\subsection{维特定理。}

给定 A 上的两个向量空间 E、E'，分别配备两个半双线性型 \(\Phi\) 、\(\Phi'\)（相应地，配备两个二次型 Q、Q'），我们把从 E 到 E' 的线性映射 \(u\) 称为从 E 到 E' 的度量同态，如果它满足 \(\Phi'(u(x), u(y)) = \Phi(x, y)\)（相应地，Q'(u(x)) = Q(x)），其中 \(x \in E\) 、\(y \in E\) 。若 E 与 E' 有相同的有限维数且 \(\Phi\)（相应地，Q）非退化，则从 E 到 E' 的每个度量同态 \(u\) 都是同构，因为 \(u(x) = 0\) 蕴含 \(\Phi(x, y) = 0\) （对一切 \(y \in E\) ），从而 \(x = 0\) ；于是 \(u\) 是单射，从而因 E 与 E' 有相同的有限维数而是双射。

定理 1（维特）. — 设 E、E' 是两个有限维向量空间，分别配备满足第 2 小节条件 (T) 的两个非退化 ε-埃尔米特型 Φ 与 Φ'（相应地，配备两个非退化二次型 Q 与 Q'），并且关于这些结构彼此同构。给定 E 的任一子空间 F，F 到 E' 中的每个单射度量同态都可扩张为从 E 到 E' 上的度量同构。

利用给定的从 E 到 \(\mathbf{E}'\) 上的同构，可以看出只需证明 F 到 E 中的每个单射度量同态 \(u\) 都可扩张为 E 的度量自同构。注意，若对 \(i = 1, 2\) ，\(\mathbf{F}_i\) 是 E 的子空间，而 \(u_i\) 是 \(\mathbf{F}_i\) 到 E 中的度量同态，使得 \(\mathbf{F}_1 \cap \mathbf{F}_2 = \{0\}\) ，且 \(\Phi(u_1(x_1), u_2(x_2)) = \Phi(x_1, x_2)\) 对 \(x_i \in \mathbf{F}_i\) （ \(i = 1, 2\) ）成立，则同态 \(\nu : x_1 + x_2 \to u_1(x_1) + u_2(x_2)\) （从 \(\mathbf{F}_1 + \mathbf{F}_2\) 到 E 中，扩张 \(u_1\) 与 \(u_2\) ）是度量同态：实际上，对任取的 \(x_i, y_i\) （属于 \(\mathbf{F}_i\) ， \(i = 1, 2\) ），表达式 \(\Phi(x_1 + x_2, y_1 + y_2)\) 与 \(\Phi(u_1(x_1) + u_2(x_2), u_1(y_1) + u_2(y_2))\) （相应地，\(\mathrm{Q}(x_1 + x_2)\) 与 \(\mathrm{Q}(u_1(x_1) + u_2(x_2))\) ）中每一个的展开式都按所作假设包含两两相等的四项（相应地，三项）。此外，若 \(u_1\) 与 \(u_2\) 都是单射且 \(u_1(\mathbf{F}_1) \cap u_2(\mathbf{F}_2) = \{0\}\) ，则 \(\nu\) 是单射。

\begin{enumerate}
\def\labelenumi{\arabic{enumi})}
\tightlist
\item
  先在 u 的不动点集是 F 的超平面 U 的情形下证明维特定理。此时形如 \(u(x) - x\) （其中 \(x \in F\) ）的向量所成的集是一条直线 D。若 \(F'\) 是与 D 正交的子空间，使得 \(F' \cap F = F' \cap u(F) = \{0\}\) ，则将有 \(\Phi(u(x), y) = \Phi(x, y)\) （对 \(x \in F\) 与 \(y \in F'\) ）；因此开头的注记适用于 u 及 \(F'\) 到 E 中的恒等映射，它表明 u 可扩张到 \(F + F'\) 并且保持 \(F'\) 的点不动；形如 \(u(x) - x (x \in F + F')\) 的向量所成的集仍是直线 D。现在，对 \(x \in F, y \in F\) 我们有
\end{enumerate}

\[
\Phi (u (x), u (y) - y) = \Phi (u (x), u (y)) - \Phi (u (x), y) = \Phi (x - u (x), y),
\]

当 \(x \in \mathrm{U}\) （即 \(u(x) = x\) ）时，上式表明 \(x \in \mathrm{D}^0\) ；换言之，我们有 \(\mathrm{U} \subset \mathrm{D}^0\) 。我们区分两种情形：

\begin{enumerate}
\def\labelenumi{\alph{enumi})}
\item
  \(F \subset D^{0}\) 。公式 (5) 表明 \(u(F)\) 不包含于 \(D^{0}\) ，因而 \(F \cap D^{0} = u(F) \cap D^{0} = U\) 。于是可以取 \(F'\) 为 U 在 \(D^{0}\) 中的一个补；由于 \(F + F'\) 包含超平面 \(D^{0}\) 且不同于它，我们有 \(F + F' = E\) ，于是我们在这种情形找到了 u 到 E 的所求扩张。
\item
  \(\mathbf{F} \subset \mathbf{D}^0\) 。公式 (5) 表明 \(u(\mathbf{F}) \subset \mathbf{D}^0\) ，从而
\end{enumerate}

D ⊂ D⁰ ；于是直线 D 是迷向的（相应地，奇异的，因为对 x ∈ F 有 Q(u(x) - x) = Q(u(x)) - Φ(x, u(x)) + Q(x) = 2Q(x) - Φ(x, x) = 0）。在这些条件下，我们将证明存在 D⁰ 的一个子空间 F'，它同时是 F 与 u(F) 在 D⁰ 中的补。若 F = u(F)，这是直接的。否则，取向量 x 与 y 使得 x ∈ F，x ∉ U，y ∈ u(F)，y ∉ U；此时有 F = U + Ax，u(F) = U + Ay，且 F 不包含 x + y，否则 y = (x + y) - x 将属于 F ∩ u(F) = U；同样可以看出 x + y 不属于 u(F)；于是直线 A(x + y) 是 F 与 u(F) 在子空间 F + u(F) 中的一个补；随后只需令 F' = A(x + y) + G，其中 G 是 F + u(F) 在 D⁰ 中的一个补。这样一来，我们有 F + F' = u(F) + F' = D⁰，并且在这种情形，1) 开头所说者表明存在 u 到 E 的超平面 D⁰ 的一个扩张，并且 D⁰ 对于这个扩张是稳定的。

于是我们被归结到 F 是超平面 D⁰ 且 u 是 F 的自同构的情形。我们来证明，对每个 z ∈ E，存在 z' ∈ E 使得

\[
\Phi (u (x), z ^ {\prime}) = \Phi (x, z)\tag{6}
\]

对一切 \(x \in \mathbf{F}\) 成立；实际上，线性形式 \(x \to \Phi(u^{-1}(x), z)\) 在 \(\mathbf{F}\) 上是 \(\mathbf{E}\) 上一个线性形式的限制，这个形式属于 \(x \to \Phi(x, z')\) 这一类型，因为 \(\Phi\) 非退化；因而 (6) 成立。此外，若 \(z \notin \mathbf{F}\) ，则存在一个向量 \(z' \in \mathbf{E}\) 满足 (6) 且使得 \(\Phi(z', z') = \Phi(z, z)\) （相应地，\(Q(z') = Q(z)\) ）。实际上，当把 \(z'\) 加上一个形如 \(u(y) - y\) （ \(y \in \mathbf{F}\) ）的 \(\mathbf{D}\) 中元素时，公式 (6) 仍成立，因为 \(\mathbf{F} = \mathbf{D}^0\) ；而由于 \(z\) 与 \(\mathbf{D}\) 不正交，第 2 小节的引理 1 使我们得出结论。于是开头的注记表明，存在一个度量同态 \(v\) （从 \(\mathbf{F} + A z = E\) 到 \(E\) 中），它扩张 \(u\) 并且把 \(z\) 变换为 \(z'\) 。由于 \(\Phi\) 非退化，\(v\) 就是所寻求的 \(E\) 的度量自同构。

\begin{enumerate}
\def\labelenumi{\arabic{enumi})}
\setcounter{enumi}{1}
\tightlist
\item
  在一般情形，我们对 \(r = \dim F\) 作归纳。\(r = 0\) 的情形是平凡的。设 \(r > 0\) ，即 \(F \neq \{0\}\) ，并设 U 是 F 的一个超平面。限制 \(u_0\) （把 \(u\) 限制到 U 上所得）依归纳假设可扩张为 E 的度量自同构 \(v_0\) 。若 \(v_0\) 扩张了 \(u\) ，定理即得证明。否则，U 是 \(v_{0}^{-1}u\) 的不变元所成的集，且由 1) 存在 E 的度量自同构 \(v_{1}\) 扩张 \(v_{0}^{-1}u\) 。于是自同构 \(v_{0}v_{1}\) 就是 u 的所求扩张。证毕。
\end{enumerate}

推论 1. --- 设对 \(i=1,2\) ，\(\mathbf{E}_{i}\) 是有限维向量空间，\(\Phi_{i}\) 是非退化的 \(\varepsilon\) -埃尔米特型，定义在 \(\mathbf{E}_{i}\) 上且满足 (T)（相应地，\(\mathbf{Q}_{i}\) 是 \(\mathbf{E}_{i}\) 上的非退化二次型），\(\mathbf{E}_{i}^{\prime}\) 与 \(\mathbf{E}_{i}^{\prime\prime}\) 是 \(\mathbf{E}_{i}\) 的两个正交子空间，且 \(\mathbf{E}_{i}\) 是它们的直和。若型 \(\Phi_{1}\) 与 \(\Phi_{2}\) （相应地，\(\mathbf{Q}_{1}\) 与 \(\mathbf{Q}_{2}\) ）等价，且它们在 \(\mathbf{E}_{1}^{\prime}\) 与 \(\mathbf{E}_{2}^{\prime}\) 上的限制等价，则它们在 \(\mathbf{E}_{1}^{\prime\prime}\) 与 \(\mathbf{E}_{2}^{\prime\prime}\) 上的限制也等价。

实际上，设 \(u\) 是从 \(\mathbf{E}_{1}^{\prime}\) 到 \(\mathbf{E}_{2}^{\prime}\) 上的度量同构。由定理 1，\(u\) 可扩张为一个度量同构 \(\nu\) （从 \(\mathbf{E}_{1}\) 到 \(\mathbf{E}_{2}\) 上）。由于 \(\Phi_{i}\) 非退化，\(\mathbf{E}_{i}^{\prime\prime}\) 是 \(\mathbf{E}_{i}^{\prime}\) 在 \(\mathbf{E}_{i}\) 中的正交补，因而 \(\nu\) 把 \(\mathbf{E}_{1}^{\prime\prime}\) 映到 \(\mathbf{E}_{2}^{\prime\prime}\) 上。证毕。

推论 2. --- 在定理 1 的假设下，E 的度量自同构群传递地作用于 E 的具有给定维数的全迷向（相应地，全奇异）子空间上。此外，若 F 是 E 的全迷向（相应地，全奇异）子空间，则从 F 到 F 上的每个双射线性映射都由 E 的一个度量自同构诱导。

推论 3. --- 设 Q 是代数闭体 A 上有限维向量空间 E 上的非退化二次型。E 的度量自同构群传递地作用于 E 的具有给定维数的非迷向子空间上。

这立即由定理 1 与命题 3 的推论 2 得出。

习题. -- 1) a) 设 K 是特征为 2 的体，J : \(\xi \to \bar{\xi}\) 是 K 的一个对合反自同构，Z 是 K 的中心。证明：若 J 在 Z 上的限制不是恒等映射，则 K 中每个元素 \(\mu\) 都在 \(\bar{\mu} = \mu\) 时具有 \(\lambda + \bar{\lambda}\) 的形式（注意 Z 中存在一个元素 \(\rho \neq 0\) ，它可以写成 \(\zeta + \bar{\zeta}\) 的形式，其中 \(\zeta \in \mathbb{Z}\) ）；于是 K 上的向量空间上的每个埃尔米特型都满足条件 (T)。

\begin{enumerate}
\def\labelenumi{\alph{enumi})}
\setcounter{enumi}{1}
\tightlist
\item
  举出特征 2 的一些体的例子，它们容许一个异于恒等映射的对合反自同构 \(\xi \to \overline{\xi}\) ，并且对它们存在元素 \(\mu = \bar{\mu}\) 不具有 \(\lambda + \bar{\lambda}\) 的形式（参见第 VIII 章，§ 11，习题 4）。
\end{enumerate}

\begin{enumerate}
\def\labelenumi{\arabic{enumi})}
\setcounter{enumi}{1}
\tightlist
\item
  设 A 是一个体，E 是 A 上的向量空间，\(\Phi\) （相应地，Q）是 E 上满足条件 (T) 的非退化 \(\varepsilon\) -埃尔米特型（相应地，E 上的非退化二次型，此时 \(\Phi\) 表示与 Q 相伴的对称双线性型）。
\end{enumerate}

\begin{enumerate}
\def\labelenumi{\alph{enumi})}
\item
  证明：要使平面 P ⊂ E 是迷向的（相应地，奇异的）且非全迷向的（相应地，非全奇异的），必须且只需它只包含一条迷向（相应地，奇异）直线（参见习题 14 e））。
\item
  假设 dim \(E \geqslant 3\) ，且 E 中存在迷向向量 \(\neq 0\) 。证明：若 P 是 E 中一个非全迷向的平面，则存在一个维数为 3 的非迷向子空间 \(V \subset E\) ，它包含迷向向量 \(\neq 0\) ，且使得 \(P \subset V\) 。
\end{enumerate}

\begin{enumerate}
\def\labelenumi{\arabic{enumi})}
\setcounter{enumi}{2}
\item
  在与习题 2 相同的假设下，证明：若 dim E ≥ 3，则 E 中每条迷向直线都是两个非迷向平面的交。
\item
  假设与习题 2 相同，并进一步假设 E 是有限维的。
\end{enumerate}

\begin{enumerate}
\def\labelenumi{\alph{enumi})}
\item
  若 \(\Phi\) （相应地，Q）的指数 v \(\geqslant 1\) ，证明：对 E 中每个迷向（相应地，奇异）向量 \(a \neq 0\) ，存在由迷向（相应地，奇异）向量组成的 E 的基 \((e_i)\) ，使得 \(e_1 = a\) （参见习题 14 e））。
\item
  设 V、W 是两个相同维数 \(r \leqslant v\) 的全迷向（相应地，全奇异）子空间；证明：存在两个极大的全迷向（相应地，全奇异）子空间 \(V_{1}\) 、\(W_{1}\) ，使得 \(V \subset V_{1}\) ，\(W \subset W_{1}\) ，且 \(V_{1} \cap W_{1} = V \cap W\) 。（若 U = V ∩ W，在 \(U^{0}/U\) 中讨论。）
\item
  设 V、W、\(V_{1}\) 、\(W_{1}\) 是四个相同维数的全迷向（相应地，全奇异）子空间，使得 \(V + W\) 与 \(V_{1} + W_{1}\) 都是非迷向的。证明：存在 E 的度量自同构 u，使得 \(u(V) = V_{1}\) 且 \(u(W) = W_{1}\) 。
\item
  设 \(f\) 是 E 上的线性形式，\(\alpha\) 是 A 中形如 \(\lambda + \varepsilon\bar{\lambda}\) 的元素（相应地，A 的一个元素）。我们考虑 E 上的半双线性型
\end{enumerate}

\[
(x, y) \rightarrow \Phi_ {1} (x, y) = \Phi (x, y) + f (x) \alpha \overline {{f (y)}}
\]

（相应地，二次型

\[
x \rightarrow \mathrm{Q} _ {1} (x) = \mathrm{Q} (x) + \alpha (f (x)) ^ {2}).
\]

证明：若 \(\Phi_{1}\) （相应地，\(Q_{1}\) ）非退化，且以 \(\nu_{1}\) 表示其指数，则我们有 \(|\nu_{1} - \nu| \leqslant 1\) 。

\begin{enumerate}
\def\labelenumi{\arabic{enumi})}
\setcounter{enumi}{4}
\tightlist
\item
  \begin{enumerate}
  \def\labelenumii{\alph{enumii})}
  \tightlist
  \item
    设 B 是一个环，\(\xi \to \overline{\xi}\) 是 B 的一个对合反自同构，\(\varepsilon\) 是 B 的中心中使 \(\varepsilon \varepsilon = 1\) 的元素。证明：若 \(\beta\) 是 B 中使 \(\beta + \varepsilon \overline{\beta} \neq 0\) 的可逆元素，则 B 中存在可逆元素 \(\mu \neq 1\) ，使得 \(\mu (\beta + \varepsilon \overline{\beta}) \mu = \beta + \varepsilon \overline{\beta}\) （证明可取 \(\mu\) 使得 \(\mu \beta \overline{\mu} = \beta\) ）。
  \end{enumerate}
\end{enumerate}

\begin{enumerate}
\def\labelenumi{\alph{enumi})}
\setcounter{enumi}{1}
\tightlist
\item
  设 A 是一个体，E 是 A 上的向量空间，\(\Phi\) 是 E 上的非退化 ε-埃尔米特半双线性型，满足 (T)。证明：若 Φ 不是交错的，则对 E 的每个非迷向超平面 H，存在 E 的一个异于恒等映射的度量自同构，它保持 H 的每个元素不变（利用 a））。
\end{enumerate}

\begin{enumerate}
\def\labelenumi{\arabic{enumi})}
\setcounter{enumi}{5}
\item
  在习题 2 的假设下，设 \(a\) 是 E 中的迷向向量 \(\neq 0\) （相应地，非奇异的迷向向量（注意这样的向量只在 A 的特征为 2 时才存在））。设 \(\lambda \in \mathbf{A}\) 满足 \(\lambda + \varepsilon \bar{\lambda} = 0\) （相应地，\(\lambda = (\mathrm{Q}(a))^{-1}\) ）；证明平延 \(x \to x + \Phi(x, a) \lambda a\) （第 II 章，§ 6，习题 7）是 E 的度量自同构；反之亦然。
\item
  在习题 2 的假设下，设 G 是 E 的度量自同构群。证明：与 G 的每个元素都可交换的 E 到其自身上的半线性双射只有 E 的位似，但以下三种情形除外：dim E = 2，G 是对应于 E 上指数为 1 的二次型的度量自同构群，且 A 是三个体 \(\mathbf{F}_{2}, \mathbf{F}_{3}\) 或 \(\mathbf{F}_{4}\) 之一（利用习题 5、6 与 3；对维数为 2 的向量空间上的二次型的情形分别讨论）。
\end{enumerate}

*8) 设 A 是一个体，E 是 A 上有限维数 \textgreater0 的向量空间，Φ 是 E 上的非退化 ε-埃尔米特半双线性型，满足 (T)。以 M(Φ) 表示 E 对于 Φ 的相似变换的乘子群（§ 6，第 5 小节）。

\begin{enumerate}
\def\labelenumi{\alph{enumi})}
\item
  设 \(V_{1}\) 、\(V_{2}\) 是 E 的两个向量子空间，维数相同，又设 \(\Phi_{1}\) 、\(\Phi_{2}\) 分别是 \(\Phi\) 在 \(V_{1}\) 、\(V_{2}\) 上的限制。要存在相似变换 u 使得 \(u(V_{1}) = V_{2}\) ，必须且只需存在 \(\alpha \in M(\Phi)\) 使得 \(\Phi_{2}\) 等价于 \(\alpha\Phi_{1}\) （利用维特定理）。
\item
  设 (F, F', G) 是 E 的一个维特分解（第 2 小节），且设 \(\Phi_0\) 是 \(\Phi\) 在非迷向子空间 G 上的限制。证明：我们有 M( \(\Phi\) ) = M( \(\Phi_0\) ) （当 G ≠ \{0\} 时）。（利用维特定理及第 2 小节的命题 2）。
\item
  证明：若 Φ 的指数 v 使得 dim E = 2v，则 M(Φ) 是 A 的中心中使 ζ̄ = ζ 的元素 ζ ≠ 0 所成的群。若 dim E = 2v + 1，则 M(Φ) 是形如 ρρ̄ 的元素所成的群，其中 ρ 跑遍 A 的中心中元素 ≠ 0 所成的乘法群（利用维特定理）。（参见 § 10，习题 18）。*
\end{enumerate}

*9) 设 A 是一个域，E 是 A 上的向量空间，Q 是 E 上的非退化二次型。称 E 的自同构 \(u\) 为关于 Q 的相似变换，如果存在 A 中元素 \(\alpha \neq 0\) ，使得 \(Q(u(x)) = \alpha Q(x)\) 对每个 \(x \in E\) 成立；此时 \(u\) 也是关于与 Q 相伴的双线性型的相似变换。在 E 的维数有限且 \(>0\) 的假设下，叙述并证明关于 Q 的相似变换的与习题 8 的结果类似的命题。*

*10) 设 A 是一个体，E 是 A 上维数 \(\geqslant 2\) 的向量空间，\(\Phi_1\) （相应地，\(\Phi_2\) ）是 E 上的非退化 \(\varepsilon_1\) -埃尔米特（相应地，\(\varepsilon_2\) -埃尔米特）半双线性型，关于 A 的一个对合反自同构 \(J_1\) （相应地，\(J_2\) ），且满足条件 (T)。证明：若 E 对于 \(\Phi_1\) 的度量自同构群是对于 \(\Phi_2\) 的相似变换群的子群，则存在 \(\alpha \in A\) 使得 \(\Phi_2 = \Phi_1\alpha\) （利用习题 5 b) 与 6）。

当假设 A 交换，并且把陈述中的 \(\Phi_{1}\) 与 \(\Phi_{2}\) 换成 E 上两个非退化二次型 \(Q_{1}, Q_{2}\) 时，证明类似的性质。*

\begin{enumerate}
\def\labelenumi{\arabic{enumi})}
\setcounter{enumi}{10}
\tightlist
\item
  设 A 是一个环，J 是 A 的一个对合反自同构，E 是一个具有有限基 \((e_{i})\) 的 A-模，\(\Phi\) 是 E 上的非退化 \(\varepsilon\) -埃尔米特型，R 是 \(\Phi\) 关于 \((e_{i})\) 的矩阵；\(\Phi\) 的度量自同构群等同于满足 \(^{t}U.R.U^{j}=R\) 的可逆矩阵 U 所成的群 G。
\end{enumerate}

\begin{enumerate}
\def\labelenumi{\alph{enumi})}
\tightlist
\item
  假设存在一个矩阵 \(P\) 使得 \(R = {}^t P + \varepsilon P^J\) 。证明：对每个矩阵 \(S\) ，只要 \({}^t S + \varepsilon S^J = 0\) 且 \(P + S\) 可逆，矩阵 \(U = (^{t}P^{J} - \varepsilon^{J}S)^{-1}(P + S)\) 就属于 G，且 \(\varepsilon I + U\) 可逆。反之（证明：对每个矩阵 \(U \in G\) ，若 \(\varepsilon I + U\) 可逆，则我们有
\end{enumerate}

\[
\varepsilon (\varepsilon I + ^ {t} U) ^ {- 1} R + \varepsilon^ {\mathrm{J}} R (\varepsilon^ {\mathrm{J}} I + U ^ {\mathrm{J}}) ^ {- 1} = R).
\]

\begin{enumerate}
\def\labelenumi{\alph{enumi})}
\setcounter{enumi}{1}
\tightlist
\item
  证明：当 \(\Phi\) 满足条件 (T) 时，a) 的条件成立。A 中方程 \(2\xi = \alpha\) 对每个 \(\alpha \in \mathbf{A}\) 有且仅有一个解的情形。
\end{enumerate}

¶ 12) 设 A 是一个域，E 是 A 上有限维数为 n 的向量空间，Φ 是 E 上的非退化 ε-埃尔米特半双线性型。

\begin{enumerate}
\def\labelenumi{\alph{enumi})}
\item
  设 \(u\) 是 E 的一个自同态；证明：若 \(r_i\) （ \(1 \leqslant i \leqslant m\) ）是 \(u\) 的相似不变量（第 VII 章，§ 5，第 1 小节，定义 1），则伴随 \(u^*\) （即 \(u\) 关于 \(\Phi\) 的伴随）的相似不变量是多项式 \(\bar{r}_i\) （ \(1 \leqslant i \leqslant m\) ），其中 \(\bar{r}_i\) 是对 \(r_i\) 的每个系数施加自同构 J 而得到的（参见第 VII 章，§ 5，习题 2）。对每个首一不可约多项式 \(p \in \mathrm{A}[\mathrm{X}]\) ，它整除 \(u\) 的极小多项式，以 \(\mathrm{F}_k(u, p)\) 表示 \((p(u))^k\) 在 E 中的核，并以 \(\mathrm{F}(u, p)\) 表示诸 \(\mathrm{F}_k(u, p)\) 对一切整数 \(k > 0\) 的并。证明：若 \(p\) 与 \(q\) 是整除 \(u\) 的极小多项式的两个不同的首一不可约多项式，则子空间 \(\mathrm{F}(u, p)\) 与 \(\mathrm{F}(u^*, \bar{q})\) 正交（利用贝祖恒等式）。最后，若 G 是 E 的向量子空间且 \(u(\mathrm{G}) \subset \mathrm{G}\) ，则有 \(u^*(\mathrm{G}^0) \subset \mathrm{G}^0\) 。
\item
  我们假设 \(uu^{*} = u^{*}u\) （即称 \(u\) 为关于 \(\Phi\) 的正规自同态的情形；参见 § 7，第 3 小节）；证明这时有 \(u^{*}(F_{k}(u,p)) \subset F_{k}(u,p)\) （对一切 \(k\) ），从而 \(u^{*}(F(u,p)) \subset F(u,p)\) 。若令 \(G(p,\bar{q}) = F(u,p) \cap F(u^{*},\bar{q})\) ，证明 E 是诸子空间 \(G(p,\bar{q})\) 的直和，且 \(G(p,\bar{q})\) 与 \(G(p_{1},\bar{q}_{1})\) 在 \(p \neq q_{1}\) 或 \(p_{1} \neq q\) 时正交；特别地，\(G(p,\bar{q})\) 在 \(p \neq q\) 时是全迷向的。证明 \(G(p,\bar{p})\) 或化为 0 或是非迷向的，且当 \(p \neq q\) 时，\(G(p,\bar{q})\) 中没有非零向量与 \(G(q,\bar{p})\) 正交（利用 \(\Phi\) 非退化这一事实）；由此推出，当 \(p \neq q\) 时，\(G(p,\bar{q})\) 与 \(G(q,\bar{p})\) 是相同维数的全迷向子空间，且 \(G(p,\bar{q}) + G(q,\bar{p})\) 是非迷向的。
\item
  假设 J 不是恒等映射，或 A 的特征不是 2，且 \(u^{*} = u\) 。设 M 是满足 \(\mathbf{M} \subset \mathrm{G}(p, \overline{p})\) 且在 \(u\) 下稳定的非迷向子空间所成的集（因而它们是 A{[}X{]}- 模 \(E_{u}\) 的子模（第 VII 章，§ 5，第 1 小节））。证明：若 M 是 M 中的极小元素，则 M 是 \(E_{u}\) 的一个不可分解子模（第 VII 章，§ 4，第 7 小节）。（假设 M 是一个不可分解子模 \(M_{1}\) 与一个子模 \(M_{2} \neq \{0\}\) 的直和，极小多项式 \(p^{h}\) 与 \(p^{k}\) （分别为 u 在 \(M_{1}\) 与 \(M_{2}\) 上的限制的）满足 \(h \geqslant k\) 。注意此时 \(M_{1}\) 必然是迷向的，且每个 \(z \neq 0\) （属于 \(M_{1}\) 且使 \(p(u).z = 0\) ）都与 \(M_{1}\) 正交（利用 \(M_{1}\) 的每个子模都是单生成的这一事实）；写出 \(z = (p(u))^{h-1}.x\) 以及 z 与 \(M_{2}\) 不正交，并由此推出必然有 k = h。然后证明存在一个不可分解子模 \(N_{2}\) （属于 \(M_{2}\) ），使得 \(p^{h}\) 是 u 在 \(N_{2}\) 上的限制的极小多项式，且 \(M_{1} + N_{2}\) 非迷向；由此得出 \(M_{2} = N_{2}\) 。最后，若 \(y \in M_{2}\) 与 z 不正交，考虑 M 中由 \(w = x + \lambda y\) 生成的子模 P，其中 \(\lambda \in A\) ，并证明可取 \(\lambda\) 使 P 非迷向，办法是证明 \(\Phi((p(u))^{h-1}.w,w) \neq 0\) ；这就导致矛盾。）
\end{enumerate}

\(d)\) 由 \(c)\) 推出 \(\mathrm{G}(p,\bar{p})\) 是两两正交的不可分解子模 \(\mathrm{H}_{i}\) 的直和。若 \(p^{h}\) 是 \(u\) 在 \(\mathrm{H}_{i}\) 上的限制的极小多项式，且 \(d\) 是 \(p\) 的次数，证明在 \(\mathrm{H}_{i}\) 中存在一个维数为 \(d\) . {[}h/2{]} 的全迷向子空间。E 不含迷向向量 \(\neq 0\) 的情形（参见 § 7，第 3 小节）。

\begin{enumerate}
\def\labelenumi{\alph{enumi})}
\setcounter{enumi}{4}
\item
  当 \(u^{*} = -u\) 或 \(u^{*}u = 1\) 时，叙述并证明与 c) 和 d) 的性质类似的性质。
\item
  举出一个例子，其中 n = 4，Φ 是对称的且指数为 2，p = \(\bar{p} = X - 1\) ，u 是正规的，E = G(p, \(\bar{p}\) )，但 E 不是 M 的极小子模的直和，并且存在 u 的一个特征向量，它不是 \(u^{*}\) 的特征向量（参见 § 7，第 3 小节）。
\end{enumerate}

\begin{enumerate}
\def\labelenumi{\arabic{enumi})}
\setcounter{enumi}{12}
\item
  假设与习题 2 的相同，并进一步假设 E 具有可数基 \((e_n)\) 。设 F 是 E 的全迷向（相应地，全奇异）子空间，且 \(F^{00} = F\) 。证明：存在一个全迷向（相应地，全奇异）子空间 \(F'\) ，使得：\(1^\circ F \cap F' = \{0\}\) ；\(2^\circ\) 存在一个基 \((a_m)_{m \in I}\) 与一个基 \((a'_m)_{m \in I}\) ，分别为 F 与 \(F'\) 的基（I 是 N 的从 0 开始的区间），使得对每对指标有 \(\Phi(a_i, a_j') = \delta_{ij}\) ；\(3^\circ (F + F')^{00} = F + F'\) 且 E 是 \(F + F'\) 与 \(G = (F + F')^0\) 的直和。（用归纳法构造一个以 E 为并的非迷向子空间的递增序列 \((L_n)\) ，使得 dim \(L_{n+1} - \dim L_n = 2\) ，并对每个 \(L_n\) 应用第 2 小节的命题 2；为构造这个序列，将对每个 \(n\) 考虑最小的整数 \(k\) ，它满足 \(e_k \notin L_n\) ，并利用 § 1 的习题 9 b)）。
\item
  设 A 是特征 2 的体，E 是 A 上有限维数为 n 的向量空间，Φ 是 E 上的非退化埃尔米特型，不一定满足条件 (T)。
\end{enumerate}

\begin{enumerate}
\def\labelenumi{\alph{enumi})}
\item
  证明：\(x \in E\) 中使 \(\Phi(x, x)\) 具有形式 \(\alpha + \bar{\alpha}\) 的那些 x 所成的集 V 是 E 的一个向量子空间。
\item
  设 \(V_{1}=V\cap V^{0}\) ，\(q=\dim V_{1}\) ，\(V_{2}\) 是 \(V_{1}\) 在 V 中的补，\(V_{3}\) 是 \(V_{1}\) 在 \(V^{0}\) 中的补。证明：存在一个基 \((e_{i})_{1\leqslant i\leqslant2q}\) （\((\mathrm{V}_{2}+\mathrm{V}_{3})^{0}=\mathrm{V}_{2}^{0}\cap\mathrm{V}_{3}^{0}\) 的），使得向量 \(e_1, \ldots, e_q\) 构成 \(V_1\) 的基，且有 \(\Phi(e_i, e_{q+j}) = \delta_{ij}\) （对 \(1 \leqslant i \leqslant q\) ，\(1 \leqslant j \leqslant q\) ）。
\item
  设 \(\mathbf{G}(\Phi)\) 是 E 的（对于 \(\Phi\) 的）度量自同构群。证明：对每个 \(u \in \mathbf{G}(\Phi)\) ，有 \(u(x) = x\) （对一切 \(x \in V^0\) ）。
\item
  对每个 \(u \in \mathbf{G}(\Phi)\) ，有 \(u(V) = V\) ；设 \(u_{v}\) 是 \(u\) 在 \(V\) 上的限制，并设 \(G_{v}\) 是由诸 \(u_{v}\) 形成的群。证明：1° 同态 \(u \to u_{v}\) （从 \(\mathbf{G}(\Phi)\) 到 \(G_{v}\) 上）的核是交换的；2° 若 \(\Phi_{2}\) 是 \(\Phi\) 在 \(V_{2}\) 上的限制，且 \(\mathbf{G}(\Phi_{2})\) 是 \(V_{2}\) 对于 \(\Phi_{2}\) 的度量自同构群，则存在一个从 \(G_{v}\) 到 \(\mathbf{G}(\Phi_{2})\) 上的同态，其核是交换的（利用 \(b\) ) 与 \(c\) )）。
\item
  我们假设 A 交换且 J 是恒等映射；设 E 是 A 上维数为 3 的向量空间，\((e_{i})_{1\leqslant i\leqslant3}\) 是 E 的一个基，\(\Phi\) 是 E 上的非退化对称型，它关于 \((e_{i})\) 的矩阵是
\end{enumerate}

\[
\left( \begin{array}{c c c} 1 & 0 & 0 \\ 0 & 0 & 1 \\ 0 & 1 & 0 \end{array} \right).
\]

证明：E 中所有迷向向量都包含在由 \(e_{2}\) 与 \(e_{3}\) 张成的超平面中（参见习题 4(a)）。举出一个只包含一条迷向直线的非迷向平面的例子（参见习题 2(a)）。证明：不存在自同构 \(u \in \mathbf{G}(\Phi)\) 使得 \(u(e_{1}) = e_{1} + e_{2}\) ，尽管有 \(\Phi(e_{1}, e_{1}) = \Phi(e_{1} + e_{2}, e_{1} + e_{2})\) 。

